\chapter{Литература}
\label{ch:literature}

Список не претендует на полноту; ориентир для читателя, который хочет сопоставить монографию с классической линией CA, обратимого вычисления и «цифровой физики''.
В тексте ссылки даны по ключу автора и года; здесь --- развёрнутый перечень по темам.

\section{Цифровая физика и дискретный космос}
\noindent
Оригинал Zuse --- \textcite{zuse1969}; для чтения на английском удобна вёрстка \textcite{zuse2012zenil}
(тот же MIT-перевод \emph{Calculating Space}, плюс послесловие German--Zenil; PDF у редактора тома).
Там же --- §\,4.1 о клеточных автоматах, «цифровых частицах'' и скорости переключения на сетке
(локальный предел сигнала vs скорость импульса в модели).
\begin{refsection}
\nocite{zuse1967,zuse1969,zuse2012zenil,fredkin1990,feynman1982,mainzer2011,tHooft2016}
\printbibliography[heading=none]
\end{refsection}

\section{Клеточные автоматы}
\noindent
Линия пространственной динамики на решётке и параллельная реализация КА-моделей
(ИВМиМГ СО РАН) --- \textcite{bandman2006,bandman2001}.
\begin{refsection}
\nocite{vonneumann1966,gardner1970,cook2004,wolfram1984,wolfram2002,vichniac1984,toffoli1987,bandman2006,bandman2001}
\printbibliography[heading=none]
\end{refsection}

\section{Обратимость и предел стирания}
\begin{refsection}
\nocite{bennett1973,landauer1961,toffoli1982}
\printbibliography[heading=none]
\end{refsection}

\section{Маделунг и гидродинамическая форма КМ}
\begin{refsection}
\nocite{madelung1927}
\printbibliography[heading=none]
\end{refsection}
